
\subsubsection{6.1 The Documentation API Gap}

The central finding is a systematic mismatch between what documentation frameworks designate as most important and what repository APIs make machine-readable. The DTS assigns its highest weights to Preprocessing (1.8) and Uses (1.5) because these sections are rare and informative when present. Our results show these sections are near-zero in structured API responses across all repositories surveyed.

This does not necessarily indicate that dataset creators fail to document preprocessing and intended uses. Rather, such documentation may exist as free-text prose in dataset cards, README files, or accompanying papers -- formats not indexed by structured API endpoints. An API-based DTS scorer will therefore systematically underestimate completeness for the highest-priority sections. Closing this gap would require either repository operators adding structured API fields for Preprocessing and Uses content, or development of free-text parsing (e.g., LLM-based extraction) to recover documentation that exists in prose form.

The proxy correlation of r = 0.989, while indicating consistent algorithmic behavior, reflects the mathematical relationship between two scoring functions that share the same binary inputs. Both weighted and unweighted scores sum the same field presence indicators with different weight vectors. The near-perfect correlation is therefore partly a consequence of this structural similarity. Human annotation (n = 120, three experts) would provide the external validity assessment that proxy correlation cannot.

\subsubsection{6.2 UCI Field Mapping Failure}

The UCI coverage of 0\% is best understood as a pipeline engineering issue rather than a finding about UCI documentation quality. The DTS scorer's configured UCI field names do not match the actual key names returned by \texttt{ucimlrepo}. This is correctable by inspecting the library's output for a small sample and updating the field mapping. Until corrected, UCI results should be excluded from cross-repository completeness comparisons, and the overall coverage rate of 91.8\% should be interpreted with the understanding that it includes 62 datasets that scored zero due to this mapping failure.

\subsubsection{6.3 Implications for Repository Design}

If repository operators were to add structured API fields for Preprocessing and Uses sections, automated documentation quality assessment could cover these high-priority areas without requiring dataset creators to produce additional documentation -- only to structure existing prose documentation into machine-readable fields. Preprocessing (weight 1.8) and Uses (weight 1.5) together account for 3.3 of the 8.0 total DTS weight (41.25\%) but contribute near-zero to current automated scores.

\subsubsection{6.4 Limitations}

\textbf{L1: Proxy validation only.} The r = 0.989 correlation measures internal consistency between two computations derived from the same binary field presence values, not agreement with human expert assessment. Human annotation of 120 datasets by three independent experts is required to establish external validity. This limitation means the pipeline's construct validity -- whether it measures what the DTS framework intends to measure -- is unverified.

\textbf{L2: Causal claims not tested.} Three planned follow-up hypotheses were not executed: whether HuggingFace's structured YAML adoption caused higher documentation completeness (H-M1), whether completeness predicts search filter eligibility (H-M2), and whether completeness predicts dataset download adoption (H-M3). This paper reports only the existence proof that automated DTS scoring is feasible, not any causal or predictive relationships.

\textbf{L3: Cross-sectional snapshot.} All data were collected in March 2026. The results represent a single time point and cannot address temporal trends in documentation quality. Longitudinal analysis would require periodic re-collection.

\textbf{L4: Binary presence scoring.} The pipeline measures whether a structured API field is populated with any non-null value, not the semantic quality, accuracy, or sufficiency of the content. A dataset with a single-word task category scores identically to one with a detailed description. DTS scores from this pipeline should be interpreted as measuring structured metadata population, not documentation quality in its full semantic sense.

\textbf{L5: Unverified citations.} Two citations central to this work have not been independently confirmed in academic databases:
\begin{itemize}
\item Rondina et al. (2025): the primary source for the DTS framework, six-section taxonomy, and all inverse-frequency weights used throughout this paper.
\item Oreamuno et al. (2024): the source for the 71.52\% undocumented HuggingFace cards statistic cited in the introduction.
\end{itemize}

Manual verification of both citations against original publications is required before any submission.

\textbf{L6: Simulation bypass in implementation.} A code-level issue (\texttt{or True} forcing synthetic annotation generation) was identified and corrected during development. While the fix was applied before final results were computed, this incident illustrates the risk of silent data simulation in automated research pipelines.

\subsubsection{6.5 Broader Impact}

This pipeline provides an automated tool for population-scale documentation audits. Automated DTS scores should be interpreted as measuring structured API field coverage -- a triage tool for identifying datasets with sparse structured metadata, not a replacement for expert assessment of dataset fitness for a given purpose.

Potential risks merit consideration. If automated DTS scores are adopted as gatekeeping criteria by repository platforms before the metric's construct validity is established through human expert comparison, datasets could be inappropriately penalized or rewarded based on metadata format rather than documentation substance. Additionally, repositories with richer structured API infrastructure (such as HuggingFace Hub) will score systematically higher than equivalent-quality datasets on repositories with sparser API schemas, which could introduce bias in cross-repository audits.
